\documentclass[12pt]{article}
\usepackage[utf8]{inputenc}
\usepackage[T1]{fontenc}
\usepackage{hyperref}
\title{Veps Language Resources for K-12 Curriculum Development}
\author{Compiled by Claude Sonnet 5 (claude-sonnet-5) for Andr\'as Kornai}
\date{}
\begin{document}
\maketitle

\section{Introduction}
Veps (vepsän kel', ISO 639-3 vep; Uralic) is spoken in Russia. It has had a Latin-script standard since 2007. The search found no modern learner textbook or course. Soviet-era school textbooks are reported to exist but were not located online.

\section{Primary Educational Resources}
No textbook with a URL was found. Omniglot notes about 30 mostly school books published in the Soviet era, and these would be hard to obtain now. Overview: \url{https://omniglot.com/writing/veps.htm}

\section{Linguistic and Reference Materials}
\begin{itemize}
\item Endangered Languages Project Veps resources. It lists a Russian--Veps dictionary and a Veps--Latvian dictionary, the latter in an older Latvian orthography. \url{https://endangeredlanguages.com/node/38618/resources}
\item Kielipankki (Language Bank of Finland) page for Veps language tools. It mentions a Veps--Finnish dictionary being developed by Heidi Niva. \url{https://www.kielipankki.fi/tools/giellalt_language_models/veps/}
\end{itemize}

\section{Religious and Community Resources}
The ELP lists New Testament texts. Kielipankki mentions the Veps newspaper \emph{Kodima}. Both are linked from the pages above.

\section{Digital and Audio Resources}
The GiellaLT finite-state morphology and proofing tools for Veps are described on the Kielipankki pages above: \url{https://www.kielipankki.fi/?p=29557}

\section{Curriculum Development Context}
Veps is taught as a subject in some schools. Public materials found are mostly dictionaries, texts and tools, so a curriculum would need new materials. Contact Veps organizations in Karelia.
\end{document}